\subsection{RQ1: To what extent can different language models generate correct code?}
\label{sec:rq1}

\newcommand{\score}[1]{%
  \ifdim #1pt < 25pt \cellcolor{red!20}#1%
  \else\ifdim #1pt < 50pt \cellcolor{orange!20}#1%
  \else\ifdim #1pt < 75pt \cellcolor{yellow!30}#1%
  \else\cellcolor{green!30}#1%
  \fi\fi\fi%
}

\newcommand{\scorecell}[2]{%
  \ifdim #1pt < 25pt \cellcolor{red!20}#2%
  \else\ifdim #1pt < 50pt \cellcolor{orange!20}#2%
  \else\ifdim #1pt < 75pt \cellcolor{yellow!30}#2%
  \else\cellcolor{green!30}#2%
  \fi\fi\fi%
}

\begin{table}[htbp]
\centering
\caption{Code generation results by model across the 563 programming tasks of EvalPlus}
\label{tab:rq1_generation}
\setlength{\tabcolsep}{3pt}
\newcommand{\lightmodelrule}{%
  \arrayrulecolor{black!15}\specialrule{0.25pt}{0.5pt}{0.5pt}\arrayrulecolor{black}%
}
\begin{tabular}{@{}>{\centering\arraybackslash}m{0.31\columnwidth}*{4}{>{\centering\arraybackslash}m{0.1475\columnwidth}}@{}}
\toprule
\small\textbf{Model} & \small\textbf{Correct} & \small\textbf{Funct.} & \small\textbf{Runtime} & \small\textbf{Syntax} \\
\midrule
\multicolumn{5}{l}{\textit{Small language models}} \\
\midrule
\shortstack[c]{Muse Glimmer\\(30B)} & \textbf{81\%} & \textbf{15\%} & \textbf{4\%} & 1\% \\
\lightmodelrule
\shortstack[c]{GPT-OSS\\(20B)} & 79\% & 17\% & 5\% & 0\% \\
\lightmodelrule
\shortstack[c]{Gemma 4\\(31B)} & 79\% & 15\% & 4\% & 3\% \\
\lightmodelrule
\shortstack[c]{Qwen 3 Coder\\(30B)} & 77\% & 18\% & 6\% & 1\% \\
\lightmodelrule
\shortstack[c]{Qwen 3.8\\(27B)} & 76\% & 19\% & 6\% & 0\% \\
\lightmodelrule
\shortstack[c]{Qwen 2.5 Coder\\(7B)} & 73\% & 21\% & 8\% & 0\% \\
\lightmodelrule
\shortstack[c]{Devstral Small 2\\(24B)} & 71\% & 24\% & 7\% & 0\% \\
\lightmodelrule
\shortstack[c]{Qwen 2.5 Coder\\(32B)} & 67\% & 25\% & 9\% & 1\% \\
\lightmodelrule
\shortstack[c]{GLM 4.7 Flash\\(30B)} & 66\% & 29\% & 7\% & 0\% \\
\lightmodelrule
\shortstack[c]{Ornith 1.5\\(35B)} & 66\% & 19\% & 7\% & 10\% \\
\lightmodelrule
\shortstack[c]{Gemma 3\\(4B)} & 62\% & 33\% & 9\% & 0\% \\
\midrule
\multicolumn{5}{l}{\textit{Large language models}} \\
\midrule
\shortstack[c]{DeepSeek V4.1\\Flash\\(552B)} & 73\% & 17\% & 6\% & 6\% \\
\lightmodelrule
\shortstack[c]{MiMo V2.5\\(310B)} & 73\% & 19\% & 7\% & 2\% \\
\bottomrule
\end{tabular}

\footnotesize{* Due to the multilabel classification, Functional and Runtime may appear more than once per sample, so the sum of columns may exceed 100\%.}
\end{table}

The first question of this study concerns the ability of language models to generate valid code, with the large models included as external comparison baselines. All evaluated models produced correct code for the majority of tasks, with correctness rates ranging from 62\% to 81\%. This represents a difference of 19 percentage points between the best and worst performing models. Muse Glimmer 30B achieved the highest correctness rate (81\%, corresponding to 456 of 563 tasks), while Gemma 3 4B achieved the lowest (62\%, corresponding to 350 tasks). GPT-OSS 20B and Gemma 4 31B also performed strongly, with 79\% correctness (445 and 443 tasks, respectively).

Regarding model size and performance, larger parameter counts did not consistently translate into higher correctness rates among the evaluated models. Although several of the best-performing models have relatively large parameter counts, model size alone does not explain the observed performance differences. A notable example is the Qwen 2.5 Coder family, where the smaller 7B model achieved 73\% correctness (413 tasks), compared with 67\% (378 tasks) for its 32B counterpart. The 7B model correctly solved 82 tasks that the 32B model failed, while the 32B model correctly solved 47 tasks that the 7B model failed. This difference was significant before correction ($p = 0.0026$), but not after applying Holm correction across all 78 pairwise comparisons ($p_{\mathrm{adj}} = 0.091$); it should therefore be interpreted as descriptive evidence rather than a confirmed pairwise difference.

The large-model baselines reinforce this weak relationship between parameter count and code-generation performance. DeepSeek V4.1 Flash and MiMo V2.5 each solved 413 of the 563 tasks (73.4\%), exactly matching the number of correct solutions produced by Qwen 2.5 Coder 7B. Their correctness was also close to Qwen 3.8 27B (76\%) and Devstral Small 2 24B (71\%), while remaining below the best SLMs, Muse Glimmer 30B (81\%), GPT-OSS 20B (79\%), and Gemma 4 31B (79\%). Thus, for these short function-generation tasks, the much larger total parameter counts of DeepSeek and MiMo did not produce a correspondingly large advantage. The overlap between the two groups indicates that SLMs and LLMs are not cleanly separated by correctness on this benchmark.

Output-contract compliance further qualifies this comparison. DeepSeek produced 33 Syntax outcomes (5.9\%), all on HumanEval tasks and all involving indentation failures; in 32 of these cases, the response omitted a function definition and returned an indented function body instead. MiMo exhibited the same pattern in 10 cases (1.8\%). By contrast, most SLMs kept Syntax outcomes between 0\% and 3\%, despite their lower parameter counts. This suggests that several SLMs understood the required output format more reliably than the substantially larger DeepSeek baseline. For this evaluation, robust instruction following and adherence to the benchmark interface were at least as important as model scale: a semantically plausible body is unusable when it does not satisfy the expected function-level response contract.

\begin{table*}[htbp]
\centering
\caption{General error detection results by model (\%). CR, FR, RR, and SR denote Correct, Functional, Runtime, and Syntax recall. NS indicates metrics computed without the Syntax class.}
\label{tab:rq2_detection}
\begin{tabular}{lrrrrrrrrr}
\toprule
\textbf{Model} & \textbf{Acc} & \textbf{MF1} & \textbf{MR} & \textbf{CR} & \textbf{FR} & \textbf{RR} & \textbf{SR} & \textbf{MF1(NS)} & \textbf{MR(NS)} \\
\midrule
\multicolumn{10}{l}{\textit{General-purpose models}} \\
\midrule
Muse Glimmer 30B & \score{73.1} & \score{53.8} & \score{54.7} & \score{83.3} & \score{55.5} & \score{27.0} & \score{52.9} & \scorecell{61.3}{61.3 (+7.5)} & \scorecell{55.3}{55.3 (+0.6)} \\
Gemma 4 31B & \score{72.4} & \score{48.5} & \score{59.1} & \score{83.7} & \score{50.1} & \score{24.1} & \score{78.4} & \scorecell{55.8}{55.8 (+7.3)} & \scorecell{52.6}{52.6 (-6.5)} \\
Qwen 3.8 27B & \score{62.6} & \score{42.6} & \score{54.0} & \score{62.9} & \score{75.9} & \score{12.7} & \score{64.7} & \scorecell{48.6}{48.6 (+5.9)} & \scorecell{50.5}{50.5 (-3.6)} \\
Qwen 2.5 Coder 7B & \score{62.5} & \score{33.7} & \score{34.8} & \score{72.3} & \score{46.7} & \score{2.5} & \score{17.6} & \scorecell{39.4}{39.4 (+5.7)} & \scorecell{40.5}{40.5 (+5.7)} \\
Qwen 2.5 Coder 32B & \score{60.8} & \score{41.0} & \score{49.3} & \score{71.1} & \score{38.1} & \score{11.4} & \score{76.5} & \scorecell{43.3}{43.3 (+2.3)} & \scorecell{40.2}{40.2 (-9.1)} \\
Qwen 3 Coder 30B & \score{60.6} & \score{37.8} & \score{40.1} & \score{63.8} & \score{65.9} & \score{9.3} & \score{21.6} & \scorecell{44.1}{44.1 (+6.3)} & \scorecell{46.3}{46.3 (+6.2)} \\
Ornith 1.5 35B & \score{56.3} & \score{35.7} & \score{44.2} & \score{55.4} & \score{78.1} & \score{8.0} & \score{35.3} & \scorecell{43.1}{43.1 (+7.4)} & \scorecell{47.2}{47.2 (+3.0)} \\
Devstral Small 2 24B & \score{53.9} & \score{38.1} & \score{42.9} & \score{54.9} & \score{63.6} & \score{17.7} & \score{35.3} & \scorecell{43.3}{43.3 (+5.2)} & \scorecell{45.4}{45.4 (+2.5)} \\
GLM 4.7 Flash 30B & \score{50.4} & \score{37.3} & \score{58.4} & \score{49.3} & \score{64.0} & \score{26.2} & \score{94.1} & \scorecell{43.5}{43.5 (+6.3)} & \scorecell{46.5}{46.5 (-11.9)} \\
Gemma 4 E2B IT & \score{48.1} & \score{33.1} & \score{33.9} & \score{47.8} & \score{65.4} & \score{16.5} & \score{5.9} & \scorecell{40.6}{40.6 (+7.5)} & \scorecell{43.2}{43.2 (+9.3)} \\
Gemma 4 E4B IT & \score{45.9} & \score{37.0} & \score{50.8} & \score{38.0} & \score{86.4} & \score{31.6} & \score{47.1} & \scorecell{41.8}{41.8 (+4.7)} & \scorecell{52.0}{52.0 (+1.2)} \\
Gemma 3 4B IT & \score{38.5} & \score{26.7} & \score{32.0} & \score{30.0} & \score{82.0} & \score{10.1} & \score{5.9} & \scorecell{32.4}{32.4 (+5.7)} & \scorecell{40.7}{40.7 (+8.7)} \\
Gemma 3 4B & \score{32.9} & \score{23.4} & \score{32.8} & \score{22.6} & \score{84.7} & \score{6.3} & \score{17.6} & \scorecell{27.3}{27.3 (+3.9)} & \scorecell{37.9}{37.9 (+5.1)} \\
GPT-OSS 20B & \score{25.1} & \score{39.0} & \score{44.7} & \score{17.5} & \score{55.8} & \score{32.9} & \score{72.5} & \scorecell{40.4}{40.4 (+1.4)} & \scorecell{35.4}{35.4 (-9.3)} \\
\midrule
\multicolumn{10}{l}{\textit{Fine-tuned models}} \\
\midrule
Best Acc & \score{77.8} & \score{50.9} & \score{48.5} & \score{85.3} & \score{69.1} & \score{27.8} & \score{11.8} & \scorecell{63.1}{63.1 (+12.2)} & \scorecell{60.8}{60.8 (+12.3)} \\
Weighted MF1 & \score{73.5} & \score{54.2} & \score{62.9} & \score{75.9} & \score{77.9} & \score{32.9} & \score{64.7} & \scorecell{62.9}{62.9 (+8.7)} & \scorecell{62.2}{62.2 (-0.6)} \\
Balanced Acc & \score{67.8} & \score{38.1} & \score{38.9} & \score{74.4} & \score{65.8} & \score{15.2} & \score{0.0} & \scorecell{50.9}{50.9 (+12.7)} & \scorecell{51.8}{51.8 (+13.0)} \\
Balanced MF1 & \score{63.7} & \score{36.0} & \score{37.1} & \score{68.6} & \score{65.8} & \score{13.9} & \score{0.0} & \scorecell{48.0}{48.0 (+12.0)} & \scorecell{49.4}{49.4 (+12.4)} \\
\midrule
\multicolumn{10}{l}{\textit{Large models}} \\
\midrule
DeepSeek v4.1 Flash & \score{75.5} & \score{51.0} & \score{47.7} & \score{89.3} & \score{51.6} & \score{14.8} & \score{35.3} & \scorecell{54.1}{54.1 (+3.1)} & \scorecell{51.9}{51.9 (+4.1)} \\
MiMo v2.5 & \score{51.5} & \score{39.7} & \score{61.1} & \score{47.2} & \score{79.7} & \score{21.5} & \score{96.1} & \scorecell{44.7}{44.7 (+5.0)} & \scorecell{49.5}{49.5 (-11.7)} \\
\bottomrule
\end{tabular}
\end{table*}

Functional Error is the predominant failure mode across all models, ranging from 15\% to 33\%. Runtime Error appears with much lower frequency, ranging from 4\% to 9\%. This pattern indicates that the main challenge is not producing Python code that cannot execute, but rather producing executable solutions that do not fully satisfy the task specification.

Syntax Errors are rare for most models (0--3\%), with Ornith 1.5 35B being a clear outlier at 10\%. However, inspection revealed that 91.2\% of Ornith's cases classified as Syntax Error correspond to failures in the expected output structure, representing 9.2\% of all its generations. The benchmark expects a complete function definition, but the model sometimes returns only the function body or another incompatible structure. According to the classification procedure adopted in this study, this belongs to the Syntax Error category, although the underlying mechanism is frequently an inconsistency in the output format rather than a syntactically malformed Python token. Gemma 4 31B also exhibited formatting issues, with 72.2\% of its Syntax Errors associated with output formation problems, impacting 2.3\% of all generations. Overall, only 20\% of the errors classified as Syntax Error were actual syntax errors, corresponding to approximately 0.29\% of all generations.

A Cochran's Q test confirmed that there are statistically significant differences in correctness rates across the 13 models ($Q = 230.97$, $p < 0.001$). Pairwise comparisons using McNemar's test with Holm correction for multiple comparisons revealed that Muse Glimmer 30B significantly outperformed most lower-ranked models, including Gemma 3 4B ($p < 0.001$), Ornith 1.5 35B ($p < 0.001$), and GLM 4.7 Flash ($p < 0.001$). The difference between Muse Glimmer 30B and GPT-OSS 20B was not statistically significant after correction, suggesting that the top-performing models form a cluster with comparable performance.

In synthesis, the evaluated models normally produce structurally executable Python code, and syntax errors are relatively rare, apart from model-specific output-contract failures. Functional failures are much more important than syntax failures, indicating that models frequently produce plausible and executable code that does not correctly implement the task specification. The large-model comparison provides no evidence of a general scale advantage on these tasks: compact models can match or outperform much larger systems when they follow the required interface more consistently.
